% shared/paginate.tex — packet-wide pagination + recto starts.
%
% Every component of a lesson is compiled as its own document, so each one
% numbers its pages from 1. After pdfunite merges them, the student packet
% would read 1 / 1 / 1 2 3 4 5 / 1 2 3 / … instead of 1 … 14.
%
% This pass rebuilds an already-merged packet so that:
%   1. page numbers run across the WHOLE document, and
%   2. every component starts on an ODD (right-hand) page — a blank verso is
%      inserted after any component that would otherwise leave the next one on
%      a verso.
% Each source page is placed at full size on a same-size sheet; a white band
% inside the bottom margin covers the component-local number, and the
% packet-wide number is printed in exactly its place.
%
% The pass is idempotent by construction: the band covers whatever number is
% already there, so running it again simply renumbers. That is what lets
% shared/unit.mk re-paginate a unit packet built from already-paginated lesson
% packets without the two numbering schemes fighting.
%
% Inserted blanks ARE numbered: the blank verso is a real page of the packet,
% so its number is printed like any other — a student flipping through sees an
% unbroken run and can tell nothing is missing.
%
% Invoked by shared/paginate.sh as:
%   xelatex -jobname=paginated '\def\PacketSource{/abs/in.pdf}%
%     \def\PacketPages{1-1,{},2-2,{},3-7,{}}% shared/paginate.tex — packet-wide pagination + recto starts.
%
% Every component of a lesson is compiled as its own document, so each one
% numbers its pages from 1. After pdfunite merges them, the student packet
% would read 1 / 1 / 1 2 3 4 5 / 1 2 3 / … instead of 1 … 14.
%
% This pass rebuilds an already-merged packet so that:
%   1. page numbers run across the WHOLE document, and
%   2. every component starts on an ODD (right-hand) page — a blank verso is
%      inserted after any component that would otherwise leave the next one on
%      a verso.
% Each source page is placed at full size on a same-size sheet; a white band
% inside the bottom margin covers the component-local number, and the
% packet-wide number is printed in exactly its place.
%
% The pass is idempotent by construction: the band covers whatever number is
% already there, so running it again simply renumbers. That is what lets
% shared/unit.mk re-paginate a unit packet built from already-paginated lesson
% packets without the two numbering schemes fighting.
%
% Inserted blanks ARE numbered: the blank verso is a real page of the packet,
% so its number is printed like any other — a student flipping through sees an
% unbroken run and can tell nothing is missing.
%
% Invoked by shared/paginate.sh as:
%   xelatex -jobname=paginated '\def\PacketSource{/abs/in.pdf}%
%     \def\PacketPages{1-1,{},2-2,{},3-7,{}}% shared/paginate.tex — packet-wide pagination + recto starts.
%
% Every component of a lesson is compiled as its own document, so each one
% numbers its pages from 1. After pdfunite merges them, the student packet
% would read 1 / 1 / 1 2 3 4 5 / 1 2 3 / … instead of 1 … 14.
%
% This pass rebuilds an already-merged packet so that:
%   1. page numbers run across the WHOLE document, and
%   2. every component starts on an ODD (right-hand) page — a blank verso is
%      inserted after any component that would otherwise leave the next one on
%      a verso.
% Each source page is placed at full size on a same-size sheet; a white band
% inside the bottom margin covers the component-local number, and the
% packet-wide number is printed in exactly its place.
%
% The pass is idempotent by construction: the band covers whatever number is
% already there, so running it again simply renumbers. That is what lets
% shared/unit.mk re-paginate a unit packet built from already-paginated lesson
% packets without the two numbering schemes fighting.
%
% Inserted blanks ARE numbered: the blank verso is a real page of the packet,
% so its number is printed like any other — a student flipping through sees an
% unbroken run and can tell nothing is missing.
%
% Invoked by shared/paginate.sh as:
%   xelatex -jobname=paginated '\def\PacketSource{/abs/in.pdf}%
%     \def\PacketPages{1-1,{},2-2,{},3-7,{}}% shared/paginate.tex — packet-wide pagination + recto starts.
%
% Every component of a lesson is compiled as its own document, so each one
% numbers its pages from 1. After pdfunite merges them, the student packet
% would read 1 / 1 / 1 2 3 4 5 / 1 2 3 / … instead of 1 … 14.
%
% This pass rebuilds an already-merged packet so that:
%   1. page numbers run across the WHOLE document, and
%   2. every component starts on an ODD (right-hand) page — a blank verso is
%      inserted after any component that would otherwise leave the next one on
%      a verso.
% Each source page is placed at full size on a same-size sheet; a white band
% inside the bottom margin covers the component-local number, and the
% packet-wide number is printed in exactly its place.
%
% The pass is idempotent by construction: the band covers whatever number is
% already there, so running it again simply renumbers. That is what lets
% shared/unit.mk re-paginate a unit packet built from already-paginated lesson
% packets without the two numbering schemes fighting.
%
% Inserted blanks ARE numbered: the blank verso is a real page of the packet,
% so its number is printed like any other — a student flipping through sees an
% unbroken run and can tell nothing is missing.
%
% Invoked by shared/paginate.sh as:
%   xelatex -jobname=paginated '\def\PacketSource{/abs/in.pdf}%
%     \def\PacketPages{1-1,{},2-2,{},3-7,{}}\input{paginate}'
% where \PacketPages is a pdfpages page list over the MERGED pdf, `{}` marking
% an inserted blank.
%
% Geometry note: the band and the baseline below are matched to the article
% preamble in apstats-article.sty (bottom margin 0.75in = 54pt, default
% \footskip = 30pt, so the footer baseline sits 24pt above the page edge).
% Measured 2026-07-30 on unit01/lesson01/notes: the footer digit's em box runs
% 22.2pt–32.1pt above the bottom edge, well inside the 10pt–42pt band, and the
% band's top stays 12pt clear of the body. If the page geometry there changes,
% re-measure with:
%   pdftotext -f <p> -l <p> -bbox packet.pdf out.html   # bbox of the footer digit

\documentclass[10pt]{article}
\usepackage[letterpaper, margin=0pt]{geometry}
\usepackage{xcolor}
\usepackage{pdfpages}
\usepackage{eso-pic}

\providecommand{\PacketSource}{}
\providecommand{\PacketPages}{-}

\setlength{\unitlength}{1pt}

% Bottom-margin band (pt above the page's bottom edge) and the baseline the
% packet-wide number sits on — the same baseline the article class uses.
\newcommand{\PgBandBottom}{10}
\newcommand{\PgBandHeight}{32pt}
\newcommand{\PgBaseline}{24.2}

\makeatletter

% Every page gets the same treatment, inserted blanks included. The white band
% is a no-op on a blank page; on a real one it hides the component's own number.
\newcommand{\packetnumber}{%
  \AddToShipoutPictureFG*{%
    \AtPageLowerLeft{%
      \put(0,\PgBandBottom){\textcolor{white}{\rule{\paperwidth}{\PgBandHeight}}}%
      \put(0,\PgBaseline){\makebox[\paperwidth]{\normalsize\thepage}}%
    }%
  }%
}

% fitpaper + noautoscale keep every source page at exactly its original size:
% no re-scaling of the component's layout, only the footer overlay is added.
\newcommand{\packetinclude}[1]{%
  \includepdf[pages={#1}, fitpaper=true, noautoscale=true,
              pagecommand=\packetnumber]{\PacketSource}}
\edef\PacketPageList{\PacketPages}

\makeatother

\pagestyle{empty}

\begin{document}
\expandafter\packetinclude\expandafter{\PacketPageList}
\end{document}
'
% where \PacketPages is a pdfpages page list over the MERGED pdf, `{}` marking
% an inserted blank.
%
% Geometry note: the band and the baseline below are matched to the article
% preamble in apstats-article.sty (bottom margin 0.75in = 54pt, default
% \footskip = 30pt, so the footer baseline sits 24pt above the page edge).
% Measured 2026-07-30 on unit01/lesson01/notes: the footer digit's em box runs
% 22.2pt–32.1pt above the bottom edge, well inside the 10pt–42pt band, and the
% band's top stays 12pt clear of the body. If the page geometry there changes,
% re-measure with:
%   pdftotext -f <p> -l <p> -bbox packet.pdf out.html   # bbox of the footer digit

\documentclass[10pt]{article}
\usepackage[letterpaper, margin=0pt]{geometry}
\usepackage{xcolor}
\usepackage{pdfpages}
\usepackage{eso-pic}

\providecommand{\PacketSource}{}
\providecommand{\PacketPages}{-}

\setlength{\unitlength}{1pt}

% Bottom-margin band (pt above the page's bottom edge) and the baseline the
% packet-wide number sits on — the same baseline the article class uses.
\newcommand{\PgBandBottom}{10}
\newcommand{\PgBandHeight}{32pt}
\newcommand{\PgBaseline}{24.2}

\makeatletter

% Every page gets the same treatment, inserted blanks included. The white band
% is a no-op on a blank page; on a real one it hides the component's own number.
\newcommand{\packetnumber}{%
  \AddToShipoutPictureFG*{%
    \AtPageLowerLeft{%
      \put(0,\PgBandBottom){\textcolor{white}{\rule{\paperwidth}{\PgBandHeight}}}%
      \put(0,\PgBaseline){\makebox[\paperwidth]{\normalsize\thepage}}%
    }%
  }%
}

% fitpaper + noautoscale keep every source page at exactly its original size:
% no re-scaling of the component's layout, only the footer overlay is added.
\newcommand{\packetinclude}[1]{%
  \includepdf[pages={#1}, fitpaper=true, noautoscale=true,
              pagecommand=\packetnumber]{\PacketSource}}
\edef\PacketPageList{\PacketPages}

\makeatother

\pagestyle{empty}

\begin{document}
\expandafter\packetinclude\expandafter{\PacketPageList}
\end{document}
'
% where \PacketPages is a pdfpages page list over the MERGED pdf, `{}` marking
% an inserted blank.
%
% Geometry note: the band and the baseline below are matched to the article
% preamble in apstats-article.sty (bottom margin 0.75in = 54pt, default
% \footskip = 30pt, so the footer baseline sits 24pt above the page edge).
% Measured 2026-07-30 on unit01/lesson01/notes: the footer digit's em box runs
% 22.2pt–32.1pt above the bottom edge, well inside the 10pt–42pt band, and the
% band's top stays 12pt clear of the body. If the page geometry there changes,
% re-measure with:
%   pdftotext -f <p> -l <p> -bbox packet.pdf out.html   # bbox of the footer digit

\documentclass[10pt]{article}
\usepackage[letterpaper, margin=0pt]{geometry}
\usepackage{xcolor}
\usepackage{pdfpages}
\usepackage{eso-pic}

\providecommand{\PacketSource}{}
\providecommand{\PacketPages}{-}

\setlength{\unitlength}{1pt}

% Bottom-margin band (pt above the page's bottom edge) and the baseline the
% packet-wide number sits on — the same baseline the article class uses.
\newcommand{\PgBandBottom}{10}
\newcommand{\PgBandHeight}{32pt}
\newcommand{\PgBaseline}{24.2}

\makeatletter

% Every page gets the same treatment, inserted blanks included. The white band
% is a no-op on a blank page; on a real one it hides the component's own number.
\newcommand{\packetnumber}{%
  \AddToShipoutPictureFG*{%
    \AtPageLowerLeft{%
      \put(0,\PgBandBottom){\textcolor{white}{\rule{\paperwidth}{\PgBandHeight}}}%
      \put(0,\PgBaseline){\makebox[\paperwidth]{\normalsize\thepage}}%
    }%
  }%
}

% fitpaper + noautoscale keep every source page at exactly its original size:
% no re-scaling of the component's layout, only the footer overlay is added.
\newcommand{\packetinclude}[1]{%
  \includepdf[pages={#1}, fitpaper=true, noautoscale=true,
              pagecommand=\packetnumber]{\PacketSource}}
\edef\PacketPageList{\PacketPages}

\makeatother

\pagestyle{empty}

\begin{document}
\expandafter\packetinclude\expandafter{\PacketPageList}
\end{document}
'
% where \PacketPages is a pdfpages page list over the MERGED pdf, `{}` marking
% an inserted blank.
%
% Geometry note: the band and the baseline below are matched to the article
% preamble in apstats-article.sty (bottom margin 0.75in = 54pt, default
% \footskip = 30pt, so the footer baseline sits 24pt above the page edge).
% Measured 2026-07-30 on unit01/lesson01/notes: the footer digit's em box runs
% 22.2pt–32.1pt above the bottom edge, well inside the 10pt–42pt band, and the
% band's top stays 12pt clear of the body. If the page geometry there changes,
% re-measure with:
%   pdftotext -f <p> -l <p> -bbox packet.pdf out.html   # bbox of the footer digit

\documentclass[10pt]{article}
\usepackage[letterpaper, margin=0pt]{geometry}
\usepackage{xcolor}
\usepackage{pdfpages}
\usepackage{eso-pic}

\providecommand{\PacketSource}{}
\providecommand{\PacketPages}{-}

\setlength{\unitlength}{1pt}

% Bottom-margin band (pt above the page's bottom edge) and the baseline the
% packet-wide number sits on — the same baseline the article class uses.
\newcommand{\PgBandBottom}{10}
\newcommand{\PgBandHeight}{32pt}
\newcommand{\PgBaseline}{24.2}

\makeatletter

% Every page gets the same treatment, inserted blanks included. The white band
% is a no-op on a blank page; on a real one it hides the component's own number.
\newcommand{\packetnumber}{%
  \AddToShipoutPictureFG*{%
    \AtPageLowerLeft{%
      \put(0,\PgBandBottom){\textcolor{white}{\rule{\paperwidth}{\PgBandHeight}}}%
      \put(0,\PgBaseline){\makebox[\paperwidth]{\normalsize\thepage}}%
    }%
  }%
}

% fitpaper + noautoscale keep every source page at exactly its original size:
% no re-scaling of the component's layout, only the footer overlay is added.
\newcommand{\packetinclude}[1]{%
  \includepdf[pages={#1}, fitpaper=true, noautoscale=true,
              pagecommand=\packetnumber]{\PacketSource}}
\edef\PacketPageList{\PacketPages}

\makeatother

\pagestyle{empty}

\begin{document}
\expandafter\packetinclude\expandafter{\PacketPageList}
\end{document}
